\begin{frame}{유효 시계: $\gamma = 0.95$ 는 몇 스텝짜리 시야인가}
  할인율을 ``몇 $\%$씩 깎는 비율''로만 이해하면, ``0.95와 0.9의
  차이가 실제로 큰가''를 느끼기 어렵다.
  \vskip0.4em
  \begin{block}{재정의}
    ``할인 계수 $\gamma^k$ 가 $\epsilon$ \textbf{아래로 떨어지는 첫
    스텝 $k$ 를 $H$ 라 하자'' --- 이 $H$ 를 \kb{유효 시계
    (effective horizon)}라고 부른다. ``에이전트는 대략 $H$ 스텝
    까지 미래를 볼 수 있다''고 읽을 수 있다.
  }\end{block}
  \vskip0.5em
  $\gamma^H = \epsilon$ 양변에 로그를 취하면:
  \[
  H = \frac{\log \epsilon}{\log \gamma}
  \]
  $\epsilon = 0.05$ 로 잡으면(= 5\% 이하로 깎인 보상부터 ``무시''):
  \small
  \begin{tabular}{@{}lll@{}}
    \toprule
    $\gamma$ & 유효 시계 $H$ ($\epsilon=0.05$) & 해석 \\
    \midrule
    0.5 & 4.3 $\to$ \textbf{약 4스텝} & 5스텝 뒤 보상은 이미 3\% \\
    0.9 & 28.4 $\to$ \textbf{약 28스텝} & 약 28스텝 뒤 보상은 5\% 이하 \\
    0.95 & 58.4 $\to$ \textbf{약 58스텝} & CartPole 에피소드(500)의 $\sim$1/9 \\
    0.99 & 298.1 $\to$ \textbf{약 298스텝} & 500스텝 에피소드의 60\%까지 보임 \\
    \bottomrule
  \end{tabular}
\end{frame}
